% Lösungen Einheit 3 von 4 – Senkrecht zu Geraden und Ebenen (zusammenbau v0.1)
\einheitenkopf{Einheit 3 von 4 · Senkrecht zu Geraden und Ebenen}

\erg{16}{a) Gerade gegen Ebene – Richtungsvektor kollinear zum Normalenvektor \quad b) Der Richtungsvektor $(2 | -1 | 3)$ von $g$ ist zugleich Normalenvektor von $E$ – die Gerade hat die Richtung der Normalen, also steht sie senkrecht auf $E$ \quad c) $\overrightarrow{BE} = (4 | -2 | 4) = 2 \cdot (2 | -1 | 2)$, ein Vielfaches des Normalenvektors von $L$ – die Seitenkanten stehen senkrecht auf der Grundfläche, das Prisma ist gerade \quad d) $(2 | 1 | 2) \circ (1 | -4 | 1) = 2 - 4 + 2 = 0$ und $(1 | 0 | -1) \circ (1 | -4 | 1) = 1 + 0 - 1 = 0$ – $\vec n$ steht auf beiden Spannvektoren senkrecht, also auf $E$ \quad e) Richtungsvektor von $h$: $(0 | 1 | 0)$; $(-3 | 0 | 2) \circ (0 | 1 | 0) = 0 + 0 + 0 = 0$ – ja, senkrecht \quad f) $(4 | -3 | 1) \circ (1 | 2 | k) = 4 - 6 + k = 0$, also $k = 2$ \quad g) z. B. $H\colon 3x - 5y - 7z = 0$ – Normalenvektor von $H$ ist der Richtungsvektor $(3 | -5 | -7)$ der Schnittgeraden (er steht auf beiden Normalenvektoren senkrecht), der Stützpunkt ist frei}
\erg{17}{a) Fehler: Paul erwartet ein Skalarprodukt null; das hieße aber, der Richtungsvektor liegt parallel zur Ebene. Senkrecht zur Ebene heißt kollinear zum Normalenvektor – hier sind beide Vektoren sogar gleich, $g$ steht senkrecht auf $E$. \quad b) Der Normalenvektor steht auf allen Richtungen der Ebene senkrecht; eine Gerade, die auf der Ebene senkrecht steht, muss ebenfalls zu allen Richtungen der Ebene senkrecht sein – das leistet nur die Richtung des Normalenvektors (bis auf ein Vielfaches). \quad c) Ja – $(2 | 0 | -8) = 2 \cdot (1 | 0 | -4)$ ist ein Vielfaches des Normalenvektors des Hangs, der Mast hat die Richtung der Normalen}
